\IfFileExists{stacks-project.cls}{%
\documentclass{stacks-project}
}{%
\documentclass{amsart}
}

% For dealing with references we use the comment environment
\usepackage{verbatim}
\newenvironment{reference}{\comment}{\endcomment}
%\newenvironment{reference}{}{}
\newenvironment{slogan}{\comment}{\endcomment}
\newenvironment{history}{\comment}{\endcomment}

% For commutative diagrams we use Xy-pic
\usepackage[all]{xy}

% We use 2cell for 2-commutative diagrams.
\xyoption{2cell}
\UseAllTwocells

% We use multicol for the list of chapters between chapters
\usepackage{multicol}

% This is generally recommended for better output
\usepackage{lmodern}
\usepackage[T1]{fontenc}

% For cross-file-references
\usepackage{xr-hyper}

% Package for hypertext links:
\usepackage{hyperref}

% For any local file, say "hello.tex" you want to link to please
% use \externaldocument[hello-]{hello}
\externaldocument[introduction-]{introduction}
\externaldocument[conventions-]{conventions}
\externaldocument[sets-]{sets}
\externaldocument[categories-]{categories}
\externaldocument[topology-]{topology}
\externaldocument[sheaves-]{sheaves}
\externaldocument[sites-]{sites}
\externaldocument[stacks-]{stacks}
\externaldocument[fields-]{fields}
\externaldocument[algebra-]{algebra}
\externaldocument[brauer-]{brauer}
\externaldocument[homology-]{homology}
\externaldocument[derived-]{derived}
\externaldocument[simplicial-]{simplicial}
\externaldocument[more-algebra-]{more-algebra}
\externaldocument[smoothing-]{smoothing}
\externaldocument[modules-]{modules}
\externaldocument[sites-modules-]{sites-modules}
\externaldocument[injectives-]{injectives}
\externaldocument[cohomology-]{cohomology}
\externaldocument[sites-cohomology-]{sites-cohomology}
\externaldocument[dga-]{dga}
\externaldocument[dpa-]{dpa}
\externaldocument[sdga-]{sdga}
\externaldocument[hypercovering-]{hypercovering}
\externaldocument[schemes-]{schemes}
\externaldocument[constructions-]{constructions}
\externaldocument[properties-]{properties}
\externaldocument[morphisms-]{morphisms}
\externaldocument[coherent-]{coherent}
\externaldocument[divisors-]{divisors}
\externaldocument[limits-]{limits}
\externaldocument[varieties-]{varieties}
\externaldocument[topologies-]{topologies}
\externaldocument[descent-]{descent}
\externaldocument[perfect-]{perfect}
\externaldocument[more-morphisms-]{more-morphisms}
\externaldocument[flat-]{flat}
\externaldocument[groupoids-]{groupoids}
\externaldocument[more-groupoids-]{more-groupoids}
\externaldocument[etale-]{etale}
\externaldocument[chow-]{chow}
\externaldocument[intersection-]{intersection}
\externaldocument[pic-]{pic}
\externaldocument[weil-]{weil}
\externaldocument[adequate-]{adequate}
\externaldocument[dualizing-]{dualizing}
\externaldocument[duality-]{duality}
\externaldocument[discriminant-]{discriminant}
\externaldocument[derham-]{derham}
\externaldocument[local-cohomology-]{local-cohomology}
\externaldocument[algebraization-]{algebraization}
\externaldocument[curves-]{curves}
\externaldocument[resolve-]{resolve}
\externaldocument[models-]{models}
\externaldocument[functors-]{functors}
\externaldocument[equiv-]{equiv}
\externaldocument[pione-]{pione}
\externaldocument[etale-cohomology-]{etale-cohomology}
\externaldocument[proetale-]{proetale}
\externaldocument[relative-cycles-]{relative-cycles}
\externaldocument[more-etale-]{more-etale}
\externaldocument[trace-]{trace}
\externaldocument[crystalline-]{crystalline}
\externaldocument[spaces-]{spaces}
\externaldocument[spaces-properties-]{spaces-properties}
\externaldocument[spaces-morphisms-]{spaces-morphisms}
\externaldocument[decent-spaces-]{decent-spaces}
\externaldocument[spaces-cohomology-]{spaces-cohomology}
\externaldocument[spaces-limits-]{spaces-limits}
\externaldocument[spaces-divisors-]{spaces-divisors}
\externaldocument[spaces-over-fields-]{spaces-over-fields}
\externaldocument[spaces-topologies-]{spaces-topologies}
\externaldocument[spaces-descent-]{spaces-descent}
\externaldocument[spaces-perfect-]{spaces-perfect}
\externaldocument[spaces-more-morphisms-]{spaces-more-morphisms}
\externaldocument[spaces-flat-]{spaces-flat}
\externaldocument[spaces-groupoids-]{spaces-groupoids}
\externaldocument[spaces-more-groupoids-]{spaces-more-groupoids}
\externaldocument[bootstrap-]{bootstrap}
\externaldocument[spaces-pushouts-]{spaces-pushouts}
\externaldocument[spaces-chow-]{spaces-chow}
\externaldocument[groupoids-quotients-]{groupoids-quotients}
\externaldocument[spaces-more-cohomology-]{spaces-more-cohomology}
\externaldocument[spaces-simplicial-]{spaces-simplicial}
\externaldocument[spaces-duality-]{spaces-duality}
\externaldocument[formal-spaces-]{formal-spaces}
\externaldocument[restricted-]{restricted}
\externaldocument[spaces-resolve-]{spaces-resolve}
\externaldocument[formal-defos-]{formal-defos}
\externaldocument[defos-]{defos}
\externaldocument[cotangent-]{cotangent}
\externaldocument[examples-defos-]{examples-defos}
\externaldocument[algebraic-]{algebraic}
\externaldocument[examples-stacks-]{examples-stacks}
\externaldocument[stacks-sheaves-]{stacks-sheaves}
\externaldocument[criteria-]{criteria}
\externaldocument[artin-]{artin}
\externaldocument[quot-]{quot}
\externaldocument[stacks-properties-]{stacks-properties}
\externaldocument[stacks-morphisms-]{stacks-morphisms}
\externaldocument[stacks-limits-]{stacks-limits}
\externaldocument[stacks-cohomology-]{stacks-cohomology}
\externaldocument[stacks-perfect-]{stacks-perfect}
\externaldocument[stacks-introduction-]{stacks-introduction}
\externaldocument[stacks-more-morphisms-]{stacks-more-morphisms}
\externaldocument[stacks-geometry-]{stacks-geometry}
\externaldocument[moduli-]{moduli}
\externaldocument[moduli-curves-]{moduli-curves}
\externaldocument[examples-]{examples}
\externaldocument[exercises-]{exercises}
\externaldocument[guide-]{guide}
\externaldocument[desirables-]{desirables}
\externaldocument[coding-]{coding}
\externaldocument[obsolete-]{obsolete}
\externaldocument[fdl-]{fdl}
\externaldocument[index-]{index}

% Theorem environments.
%
\theoremstyle{plain}
\newtheorem{theorem}[subsection]{Theorem}
\newtheorem{proposition}[subsection]{Proposition}
\newtheorem{lemma}[subsection]{Lemma}

\theoremstyle{definition}
\newtheorem{definition}[subsection]{Definition}
\newtheorem{example}[subsection]{Example}
\newtheorem{exercise}[subsection]{Exercise}
\newtheorem{situation}[subsection]{Situation}

\theoremstyle{remark}
\newtheorem{remark}[subsection]{Remark}
\newtheorem{remarks}[subsection]{Remarks}

\numberwithin{equation}{subsection}

% Macros
%
\def\lim{\mathop{\mathrm{lim}}\nolimits}
\def\colim{\mathop{\mathrm{colim}}\nolimits}
\def\Spec{\mathop{\mathrm{Spec}}}
\def\Hom{\mathop{\mathrm{Hom}}\nolimits}
\def\Ext{\mathop{\mathrm{Ext}}\nolimits}
\def\SheafHom{\mathop{\mathcal{H}\!\mathit{om}}\nolimits}
\def\SheafExt{\mathop{\mathcal{E}\!\mathit{xt}}\nolimits}
\def\Sch{\mathit{Sch}}
\def\Mor{\mathop{\mathrm{Mor}}\nolimits}
\def\Ob{\mathop{\mathrm{Ob}}\nolimits}
\def\Sh{\mathop{\mathit{Sh}}\nolimits}
\def\NL{\mathop{N\!L}\nolimits}
\def\CH{\mathop{\mathrm{CH}}\nolimits}
\def\proetale{{pro\text{-}\acute{e}tale}}
\def\etale{{\acute{e}tale}}
\def\QCoh{\mathit{QCoh}}
\def\Ker{\mathop{\mathrm{Ker}}}
\def\Im{\mathop{\mathrm{Im}}}
\def\Coker{\mathop{\mathrm{Coker}}}
\def\Coim{\mathop{\mathrm{Coim}}}

% Boxtimes
%
\DeclareMathSymbol{\boxtimes}{\mathbin}{AMSa}{"02}

%
% Macros for moduli stacks/spaces
%
\def\QCohstack{\mathcal{QC}\!\mathit{oh}}
\def\Cohstack{\mathcal{C}\!\mathit{oh}}
\def\Spacesstack{\mathcal{S}\!\mathit{paces}}
\def\Quotfunctor{\mathrm{Quot}}
\def\Hilbfunctor{\mathrm{Hilb}}
\def\Curvesstack{\mathcal{C}\!\mathit{urves}}
\def\Polarizedstack{\mathcal{P}\!\mathit{olarized}}
\def\Complexesstack{\mathcal{C}\!\mathit{omplexes}}
% \Pic is the operator that assigns to X its picard group, usage \Pic(X)
% \Picardstack_{X/B} denotes the Picard stack of X over B
% \Picardfunctor_{X/B} denotes the Picard functor of X over B
\def\Pic{\mathop{\mathrm{Pic}}\nolimits}
\def\Picardstack{\mathcal{P}\!\mathit{ic}}
\def\Picardfunctor{\mathrm{Pic}}
\def\Deformationcategory{\mathcal{D}\!\mathit{ef}}
\usepackage{cleveref}
\crefname{originallemma}{original lemma}{original lemmas}
\Crefname{originallemma}{Original lemma}{Original lemmas}
\makeatletter
\newlabel{homology-lemma-ses-termwise-split}{{011D}{source}{}{sourcepointer-1}{}}
\newlabel{homology-lemma-ses-termwise-split@cref}{{[originallemma][0][]011D}{[1][0][]source}{}{}{}}
\makeatother
\title{Splittings whose connecting map factors through (1,-1) give a map into the diamond complex}
\author{The Stacks Project Authors}
\date{stacks-project:simplicial.tex:local-snapshot-20261012; GNU FDL 1.2 or later}
\begin{document}
\maketitle
\noindent\textit{Editorial note.} The counted statement is the author's own lemma environment labelled lemma-map-into-diamond (source lines 5433-5467); the proof environment begins on the second line after it and is its own proof, closing at source line 5504.\par\medskip
\noindent\textit{Editorial note.} The counted spans invoke no labelled result of this chapter, so the card carries no context span; the one cross-chapter target (the termwise-splitting lemma of Homology) carries a canonical project tag and is auto-declared as a non-load-bearing source pointer.\par\medskip
\noindent\textit{Editorial note.} The retained preamble is the chapter's own minimal preamble (sources/stacks-simplicial/preamble.tex, taken with its real line count), a verbatim excerpt of the native Xy-pic setup, the native theorem declarations with the mathematical macros, and the project's own no-op reference environments; no macro was invented or redefined.\par\medskip
\noindent\textit{Editorial scope.} One original result of the Stacks Project's Simplicial Methods chapter, counted once: termwise splittings whose connecting map factors through (1,-1) give a unique map of the extension into the diamond complex, compatible with the splittings.\par\medskip
\noindent\textit{Allowed standard prerequisites (editorial).}\begin{itemize}
\item Chain complexes and short exact sequences. Short exact sequences of chain complexes in an abelian category, with morphisms written as components a, b, r.. It is the setting of the statement and of the computation (source lines 5435-5440, 5472-5473).
\item Termwise splittings and the connecting morphism. Given splittings s\_n of the sequence in each degree, the projections (p\_n, q\_n) and the connecting morphism delta(s) : C to A[-1] direct sum A[-1] are defined (Homology, Lemma homology-lemma-ses-termwise-split, declared as a tagged source pointer).. It fixes the notation and the meaning of the hypothesis (source lines 5441-5445, 5470-5471).
\item The diamond complex. The complex diamond A with diamond A\_n = A\_n direct sum A\_n direct sum A\_\{n-1\} and an explicit matrix differential, which represents triples of two maps and a homotopy.. It is the target of the constructed morphism and its differential is what makes the third term of the computation come out (source lines 5448-5466, 5478, 5495).
\end{itemize}\medskip
\ifdefined\hypertarget\hypertarget{sourcepointer-1}{}\fi
\noindent\textit{Source reference.} Original lemma 011D: sources/stacks-chow/tags.txt, source lines 1359--1359. Cross-chapter pointer whose target carries the canonical tag 011D in the project tag list; declared so the reference resolves inside the assembled package, while every step the counted argument uses remains inside the retained spans.\par\medskip
\begin{lemma}
\label{lemma-map-into-diamond}
Let $\mathcal{A}$ be an abelian category.
Let
$$
0 \to A \oplus A \to B \to C \to 0
$$
be a short exact sequence of chain complexes of $\mathcal{A}$.
Suppose given in addition morphisms $s_n : C_n \to B_n$
splitting the associated short exact sequence in degree $n$.
Let $\delta(s) : C \to (A \oplus A)[-1] = A[-1] \oplus A[-1]$
be the associated morphism of complexes, see
Homology, Lemma \ref{homology-lemma-ses-termwise-split}.
If $\delta(s)$ factors through the morphism
$(1, -1) : A[-1] \to A[-1] \oplus A[-1]$, then
there is a unique morphism $B \to \diamond A$
fitting into a commutative diagram
$$
\xymatrix{
0 \ar[r] &
A \oplus A \ar[d] \ar[r] &
B \ar[r] \ar[d] &
C \ar[d] \ar[r] &
0 \\
0 \ar[r] &
A \oplus A \ar[r] &
\diamond A \ar[r] &
A[-1] \ar[r] &
0
}
$$
where the vertical maps are compatible with the splittings
$s_n$ and the splittings of $\diamond A_n \to A[-1]_n$
as well.
\end{lemma}

\begin{proof}
Denote $(p_n, q_n) : B_n \to A_n \oplus A_n$ the morphism $\pi_n$
of Homology, Lemma \ref{homology-lemma-ses-termwise-split}.
Also write $(a, b) : A \oplus A \to B$, and
$r : B \to C$ for the maps in the short exact sequence.
Write the factorization of $\delta(s)$ as
$\delta(s) = (1, -1) \circ f$. This means that
$p_{n - 1} \circ d_{B, n} \circ s_n = f_n$, and
$q_{n - 1} \circ d_{B, n} \circ s_n = - f_n$, and
Set $B_n \to \diamond A_n = A_n \oplus A_n \oplus A_{n - 1}$
equal to $(p_n, q_n, f_n \circ r_n)$.

\medskip\noindent
Now we have to check that this actually defines a morphism
of complexes. We will only do this in the case of abelian groups.
Pick $x \in B_n$. Then $x = a_n(x_1) + b_n(x_2) + s_n(x_3)$
and it suffices to show that our definition commutes
with differential for each term separately. For the term
$a_n(x_1)$ we have $(p_n, q_n, f_n \circ r_n)(a_n(x_1)) =
(x_1, 0, 0)$ and the result is obvious. Similarly for
the term $b_n(x_2)$. For the term $s_n(x_3)$ we have
\begin{eqnarray*}
(p_n, q_n, f_n \circ r_n)(d_n(s_n(x_3))) & = &
(p_n, q_n, f_n \circ r_n)( \\
& & \ \ \ \ \ a_n(f_n(x_3)) - b_n(f_n(x_3)) + s_n(d_n(x_3))) \\
& = &
(f_n(x_3), -f_n(x_3), f_n(d_n(x_3)))
\end{eqnarray*}
by definition of $f_n$. And
\begin{eqnarray*}
d_n(p_n, q_n, f_n \circ r_n)(s_n(x_3)) & = & d_n(0, 0, f_n(x_3)) \\
& = &
(f_n(x_3), - f_n(x_3), d_{A[-1], n}(f_n(x_3)))
\end{eqnarray*}
The result follows as $f$ is a morphism of complexes.
\end{proof}

\end{document}
